\section{Additional Related Work}
Evaluation of LLMs for programming has evolved into repository and project-level benchmarks including SWE-Bench, RepoBench, ExecRepoBench, and ProjectEval \citep{jimenez2024swebenchlanguagemodelsresolve,liu2023repobench,yang2024execrepobench,liu2025projecteval}. More recent efforts have incorporated performance-aware metrics, measuring efficiency, runtime, or resource usage in controlled settings (e.g. EffiBench, COMPASS) \citep{huang2024effibench,meaden2025compass}. While these benchmarks represent important progress beyond toy synthesis tasks, they struggle to capture real-world software engineering behaviors that are open-ended, user-facing, and remain prone to saturation. Parallel lines of work have begun to target web development and agentic code improvement such as WebSight, Design2Code, Web2Code, FrontendBench, and Web-Bench introduce browser-based execution and UI validation \citep{bhathal2025websight,si2025design2code,yun2024web2code,zhu2025frontendbench,xu2025web}, but primarily emphasize visual fidelity or functional interactivity rather than performance as experienced by end users. 



% \subsection{Benchmarks for General Code Generation}
% \somesh{This subsection not needed anymore, frame it as evolution in a few lines in the next one}
% A large body of work has focused on evaluating LLMs for code generation using standardized benchmarks. \emph{HumanEval} \citep{chen2021evaluating} and \emph{MBPP} \citep{austin2021program} introduced small-scale programming problems in Python, where success is measured via unit test completion. These benchmarks quickly became the de facto standard for measuring code generation performance. Subsequent efforts broadened coverage: \emph{MultiPL-E} added cross-lingual tasks across multiple programming languages \citep{cassano2023multipl}, while \emph{APPS} targeted competitive programming challenges with varying levels of difficulty \citep{hendrycks2021apps}. However, these tasks remain algorithmic and fail to reflect the complexities of large-scale software engineering.

% Recent benchmarks have aimed to bridge this gap by moving from isolated problems to repository- and project-level contexts. \emph{SWE-bench} evaluates whether LLMs can resolve real-world GitHub issues end-to-end, including editing code and passing regression tests \citep{jimenez2023swebench}. \emph{RepoBench} \citep{zhou2024repobench} and \emph{ExecRepoBench} \citep{yang2025execrepobench} benchmark repository-level completion and execution, while \emph{ProjectEval} explicitly assesses multi-step programming agents at the project scale \citep{smith2025projecteval}. These works represent an important shift from toy examples toward realistic developer workflows. Nonetheless, the focus remains on general-purpose software engineering, with little attention to web development despite its ubiquity.

% \subsection{Benchmarks for Performance-Aware Code Generation}
% The evolution of code generation evaluation has progressed from a narrow focus on functional correctness toward performance-aware and multidimensional assessment, reflecting the growing sophistication of programming systems. This trajectory began with the introduction of \emph{HumanEval} in 2021 \somesh{The last subsection is not relevant now, start with this only, rather show the evolution as functional correctness (HumanEval, MBPP) -> Repository Level (RepoBench) to Performance Aware (EffiBench etc) and show how we are going from toy benchmarks to real world software engineering like SWE-bench, then go into the limitations of current SWE-Bench Effi-Bench on two axes (1) Web Performance / Development (2) Not self evolving, making them prone to contanimation etc}, which established functional correctness as the foundational benchmark through 164 hand-crafted Python problems and the pass@k evaluation methodology \cite{chen2021evaluating,runloop2025humaneval}. As models quickly advanced, moving from Codex's initial 28.8\% to current systems surpassing 96.3\% pass@1 \cite{runloop2025humaneval}, the limitations of correctness-only evaluation became increasingly apparent.

% To address this gap, \emph{Mercury} introduced the \emph{Beyond} metric, jointly capturing functional correctness and runtime efficiency \cite{park2024mercury,park2024mercury-neurips}. In parallel, \emph{COFFE} advanced evaluation by incorporating CPU instruction counts and stressful test cases to approximate time complexity \cite{li2025coffe,li2024notes}. Building further, \emph{EffiBench} expanded the scope through six distinct efficiency measures, including execution time, memory usage, and normalized performance metrics \cite{chen2024effibench,neurips2024effibench}. Most recently, comprehensive frameworks such as \emph{COMPASS} have emerged, evaluating code across correctness, efficiency, and code quality by leveraging industry-standard analysis tools and large-scale competitive programming datasets \cite{yunpeng2025compass,themoonlight2025compass}.

% \subsection{Benchmarks for Web Code Generation}
% \somesh{Merge this too, with next subsection, extend the limitations of Web Experience in last paragraph to this and the next subsection, WebBench is good, keep it because it introduces Puppeeteer etc}
% Most established code-generation benchmarks (e.g., HumanEval, MBPP) focus primarily on small algorithmic problems in languages such as Python, where correctness is measured via unit tests \citep{chen2021evaluating, austin2021program}. These general-purpose benchmarks largely ignore \emph{web-specific technologies}: there are no dedicated benchmarks for CSS, and only minimal coverage of HTML and JavaScript. Even multi-language suites such as MultiPL-E include only a limited subset of front-end tasks, without capturing the unique requirements of interactive or visual web development \citep{cassano2023multipl}.

% Recently, researchers have proposed benchmarks targeting web user interfaces. \citet{laurencon2024websight} introduced \emph{WebSight}, a dataset of 2M screenshot–HTML/CSS pairs, while \citet{si2024design2code} proposed \emph{Design2Code}, a similar image-to-code benchmark. These emphasize visual fidelity—how closely generated code matches reference screenshots—but omit dynamic behaviors and user interactions. Likewise, \citet{yun2024web2code} introduced \emph{Web2Code}, which evaluates only pixel-level similarity, not functional correctness or accessibility. Thus, these datasets remain narrow in scope and do not reflect the full complexity of web development.

% To address these gaps, more comprehensive benchmarks have recently been introduced. \emph{FrontendBench} \citep{frontendbench2025} curates 148 realistic tasks, spanning utilities, mini-games, UI components, and visualizations. It employs an automated browser-based evaluator that verifies not only DOM elements and layout, but also interactivity and event handling. Similarly, \emph{Web-Bench} \citep{webbench2025} comprises 50 multi-step projects designed by expert engineers, covering web standards (HTML, CSS, JavaScript, DOM, SVG, WebGL) and popular frameworks. Unlike unit-test-only evaluations, these benchmarks integrate browser automation tools (e.g., Playwright, Puppeteer) to perform end-to-end validation. Crucially, state-of-the-art LLMs that excel on classic benchmarks struggle on these web-centric ones—achieving only $\sim$25\% pass rates on Web-Bench tasks \citep{webbench2025}. This underscores both the difficulty and the necessity of dedicated benchmarks for web development, particularly for performance and accessibility.

% \subsection{Benchmarks for Web Performance Optimization}
% While software engineering benchmarks like SWE-bench \citep{jimenez2023swebench} and web-specific benchmarks such as Web-Bench \citep{webbench2025} and FrontendBench \citep{frontendbench2025} have advanced to project-level complexity, their primary evaluation criterion remains functional correctness (i.e., whether the code works and passes tests). These benchmarks confirm whether the code works, but not how well it performs for the user.\somesh{This makes the motivation weak, show how Performance is closer to real world engineering and why it matters more, as opposed to `just this but not this' framing}

% This omission is critical, as web quality is judged by user experience, which is quantified by Core Web Vitals (CWV) \citep{googlewebvitals}. Optimizing for CWV metrics—Largest Contentful Paint (LCP), Interaction to Next Paint (INP), and Cumulative Layout Shift (CLS) \citep{googlecwvdefinition}—requires deep, non-trivial code modifications (e.g., critical CSS inlining, image optimization) that go beyond unit-testable logic \citep{googlecwvoptimize}.

% However, no current public benchmark provides a measurable environment to test an agent's ability to propose and execute continuous, graded improvements solely against CWV scores. Existing efficiency-focused benchmarks, like Mercury \citep{du2024mercury}, have limited scope to the Python stack and only capture back-end runtime, not client-side rendering complexity. Therefore, a new benchmark is needed to rigorously assess LLMs and agentic tools for quantitative web performance optimization.

% \subsection{Benchmarks for Code Generation Frameworks}



% \subsection{Agentic Frameworks for Code Improvement}
% \somesh{Make this independent for now, and bring the code generation relevant things from agent perspective here}
% In parallel, the community has begun developing \emph{agentic frameworks} that place LLMs in iterative loops of proposing, testing, and refining code. Such systems are well-suited for error correction and quality improvements. \citet{phan2024hyperagent} introduced \emph{HyperAgent}, a multi-agent framework with roles for planning, editing, executing, and verifying code. This architecture enables systematic refinement, outperforming baselines on real-world issue resolution benchmarks such as SWE-Bench. Similarly, \citet{shinn2023reflexion} proposed the \emph{Reflexion} framework, where an agent explicitly reflects on mistakes after failures and retries with improved reasoning. Reflexion boosted code reliability dramatically, achieving 91\% pass@1 on HumanEval, surpassing GPT-4's 80\%. Other approaches such as \emph{MGDebugger} adopt hierarchical strategies: decomposing programs into components, testing each, and fixing bugs iteratively \citep{shi2025mgdebugger}.

% For web accessibility, iterative agents have also been explored. \citet{othman2023accessibility} and \citet{aljedaani2024accessible} showed that static prompting often produces inaccessible web code, while multi-turn prompting can help repair issues such as missing alt-text or ARIA labels. Building on this, \citet{lopezgil2025feeda11y} proposed \emph{FeedA11y}, a feedback-driven agent that generates code, runs accessibility checkers, and revises the output iteratively. By integrating external evaluation tools (e.g., IBM's accessibility scanner) into the generation loop, FeedA11y significantly improved WCAG compliance compared to single-shot generation.

% Our work builds on these agentic paradigms but extends them to \emph{web performance optimization}. Instead of producing a single candidate, our agent iteratively explores multiple fixes and evaluates them against metrics such as Core Web Vitals and accessibility scores. This shifts evaluation from binary unit-test outcomes to automated verification with continuous, graded signals. By allowing multiple solution attempts, varying levels of difficulty, and a self-improvement environment, our framework aims to push frontier LLMs toward more robust and optimal web development capabilities. In this respect, it unifies two trends: (1) the emergence of web-specific benchmarks like FrontendBench and Web-Bench, and (2) agentic LLM systems that learn through planning, tool-use, and feedback \citep{phan2024hyperagent, shinn2023reflexion}.